\subsection{Reading and Synthesizing Research Papers}

Choose reading depth according to the intended task: triage, comprehension, verification, reproduction, or synthesis.

\begin{process} \label{proc:three-pass-paper-reading} \hlt{Three-Pass Paper Reading}\\
Adapt the three-pass method of \citet{keshav2007readpaper} to the paper's evidence type and record the output of each completed pass.
\begin{enumerate}[label=(\roman*)]
\setlength{\itemsep}{0pt}
\item \hlt{Triage}: inspect the title, abstract, introduction, headings, conclusion, figures, and references. Record the paper's category, context, central claimed contribution, apparent assumptions, and relevance to the intended task.
\item \hlt{Comprehension}: read the body carefully, map each central claim to its stated assumptions and evidence, inspect figures and tables, and identify unfamiliar prerequisites or missing details.
\item \hlt{Reconstruction}: independently rebuild each decisive proof, derivation, algorithm, or experiment far enough to identify hidden assumptions, fragile steps, and reproducibility barriers.
\end{enumerate}
Stop after the earliest pass whose recorded output resolves the intended task.\\
If a missing prerequisite prevents the next pass, identify and acquire it before resuming.
\end{process}

\begin{process} \label{proc:paper-claim-evidence-ledger} \hlt{Paper Claim--Evidence Ledger}\\
Create one entry for each central claim and record its locator, exact scope, assumptions, evidence type, dependencies, result, uncertainty, and limitations. For a mathematical paper, reconstruct the definitions, hypotheses, cited dependencies, and proof obligations. For an empirical paper, reconstruct the sample, protocol, comparator, metric, variability, and analysis. For a systems paper, reconstruct the boundary, workload, implementation, platform, and failure model. After comparing the claim with that record, classify it as supported, partially supported, unsupported, or not assessed. Reputation and venue are metadata rather than claim-level evidence.
\end{process}

\begin{process} \label{proc:literature-synthesis} \hlt{Literature Search and Synthesis}\\
Build a claim-centered literature map rather than a sequence of paper summaries.
\begin{enumerate}[label=(\roman*)]
\setlength{\itemsep}{0pt}
\item \hlt{Seed}: start from a precise research question, corresponding keywords, a recent survey, and several directly relevant papers.
\item \hlt{Backward search}: follow repeated foundational citations to original results, data, or methods.
\item \hlt{Forward search}: locate corrections, replications, extensions, and critiques that cite each key work.
\item \hlt{Cluster}: group papers by question, assumptions, method, evidence, and operating context rather than by author.
\item \hlt{Compare}: normalize terminology and record agreements, contradictions, gaps, and incomparable protocols claim by claim.
\item \hlt{Saturate}: record the search date, sources, and inclusion boundary. Stop when another complete search and citation-tracing pass changes neither the paper clusters nor the identified agreements, contradictions, and gaps.
\end{enumerate}
Keshav's survey procedure likewise expands seed papers through repeated authors, shared references, recent proceedings \citep{keshav2007readpaper}.
\end{process}

\begin{warning} \label{warn:paper-version-and-source-status} \hlt{Version and Source Status}\\
Before treating a paper as evidence, check its version, venue, peer-review status, corrections, retractions, supplements, artifact links, and conflicts of interest. Cite the version examined and, when available, the original source for a result or dataset. Preprint and peer-review status affect provenance and scrutiny, but neither status establishes or refutes the paper's claims.
\end{warning}

\begin{principle} \label{princ:read-against-the-claim} \hlt{Read Against the Claim}\\
For each central claim, compare its scope in the title, abstract, or conclusion with the assumptions and evidence in the body. Record what was established, under which assumptions, by which evidence, with what uncertainty, and within which population, system, or operating boundary. State the synthesized conclusion no more broadly than this comparison supports.
\end{principle}
